\documentclass[12pt]{article}
\usepackage[utf8]{inputenc}
\usepackage[T1]{fontenc}
\usepackage{hyperref}
\title{Azoy\'u Me'phaa Language Resources for K-12 Curriculum Development}
\author{Compiled by Claude Sonnet 5 (claude-sonnet-5) for Andr\'as Kornai}
\date{}
\begin{document}
\maketitle

\section{Introduction}
Azoy\'u Me'phaa (Azoy\'u Tlapanec, ISO 639-3 tpc) is an Otomanguean language of Guerrero, Mexico. Search returned mostly general pages on Tlapanec (Me'phaa) as a whole and very little specific to the Azoy\'u variety. Speaker figures in the sources measure different things and conflict, so they are not repeated here.

\section{Primary Educational Resources}
No Azoy\'u-specific primers or textbooks were found in search. A UNAM-hosted document describes a 2007 Me'phaa language development congress that drew over 200 indigenous teachers: \url{https://www.nacionmulticultural.unam.mx/edespig/catalogos/noticias/documentos/doc_320.pdf}

\section{Linguistic and Reference Materials}
\begin{itemize}
\item Glottolog, Azoy\'u Tlapanec (alternate names and references): \url{https://glottolog.org/resource/languoid/id/azoy1236}
\item Wikipedia, Tlapanec language (covers the Tlapanec group as a whole): \url{https://en.wikipedia.org/wiki/Tlapanec_language}
\item Endangered Languages Alliance page on Tlapanec (says lexical resources exist and that SIL Mexico published papers on Tlapanec grammar; these are not specific to Azoy\'u): \url{https://www.elalliance.org/languages/tlapanec}
\item Rosetta Project entry for tpc: \url{https://rosettapanglossia.longnow.org/index.php?title=Tpc}
\end{itemize}

\section{Religious and Community Resources}
\begin{itemize}
\item Joshua Project profile of Azoyu Tlapaneco: \url{https://joshuaproject.net/people_groups/10529}
\end{itemize}

\section{Digital and Audio Resources}
The Rosetta Project page notes Tlapanec-language programming on radio XEZV-AM from Tlapa de Comonfort, Guerrero (see the URL above). No Azoy\'u-specific audio or corpus was found.

\section{Curriculum Development Context}
The language is not known to be taught in schools, and no dictionary or learner course for the Azoy\'u variety was found. Joshua Project says the variety may be endangered because Spanish is dominant. Materials for neighbouring Me'phaa varieties may be adaptable, but this was not verified.
\end{document}